\chapter{Auxiliary Lattice, Decision, and Fourier Constructions}
\label{app:auxiliary-mathematical-constructions}
\index{finite barrier!lattice matching}
\index{Pareto order}
\index{Bessel function}
\index{Bain path}

\chapterstatus{bounded mathematical exposition of additional repository-owned
source families. No numerical experiment, material result, new implementation,
complete source-signature crosswalk, backend admission or scientific validation
is claimed. Algebraic, analytical, software and proposed-campaign statements
below have distinct scopes.}

\section{A finite lattice barrier and its interface coordinates}

Let $j\in\mathbb Z$ label lattice sites and let $\psi_j\in\mathbb C$ be a
scalar lattice amplitude. In the source-neutral convention the negative
centered difference has coefficient one. The energy $\varepsilon$ and onsite
potential $v_j$ are dimensionless; a physical interpretation requires a
separately chosen positive energy scale $\kappa$, with physical energy
$\kappa\varepsilon$ and onsite energy $\kappa v_j$. No particular mass,
lattice spacing or semiconductor parent is selected by this normalization.
The zero-residual equation is
\begin{equation}
  2\psi_j-\psi_{j-1}-\psi_{j+1}+(v_j-\varepsilon)\psi_j=0.
  \label{eq:auxiliary-lattice-barrier-recurrence}
\end{equation}
For a constant region with potential $v$, define the two-coordinate transfer
state $s_j=(\psi_j,\psi_{j-1})^{\mathsf T}$ and transfer matrix $T(v,\varepsilon)$:
\begin{equation}
  s_{j+1}=T(v,\varepsilon)s_j,\qquad
  T(v,\varepsilon)=\begin{pmatrix}2+v-\varepsilon&-1\\1&0\end{pmatrix},
  \qquad \det T=1.
  \label{eq:auxiliary-one-site-transfer}
\end{equation}
The first row is the recurrence solved for $\psi_{j+1}$; the second row
retains $\psi_j$. A nonzero character root $z\in\mathbb C$ solves this
recurrence through $\psi_j=z^j$ exactly when
$z+z^{-1}=2+v-\varepsilon$. Its reciprocal is another root. A superposition
$\psi_j=a z^{j-o}+b z^{-(j-o)}$, with interface origin $o\in\mathbb Z$
and amplitudes $a,b\in\mathbb C$, has interface state
\begin{equation}
  s_o=S(z)\begin{pmatrix}a\\b\end{pmatrix},\qquad
  S(z)=\begin{pmatrix}1&1\\z^{-1}&z\end{pmatrix},\qquad
  \det S(z)=z-z^{-1}.
  \label{eq:auxiliary-lead-interface}
\end{equation}
Thus amplitude coordinates are nonsingular away from coincident reciprocal
roots $z=z^{-1}$, that is $z=\pm1$ for nonzero $z$.

Take a barrier of width $L\in\mathbb N$, with $v_j=v_{\mathrm L}$ for
$j<0$, $v_j=v_{\mathrm B}$ for $0\leq j<L$, and $v_j=v_{\mathrm R}$ for
$j\geq L$. Propagation over its sites gives $s_L=T(v_{\mathrm B},\varepsilon)^L s_0$.
If $z_{\mathrm L},z_{\mathrm R}$ are the respective nonzero lead roots,
matching amplitude columns $a_{\mathrm L}$ and $a_{\mathrm R}$ means
\begin{equation}
  S(z_{\mathrm R})a_{\mathrm R}
    =T(v_{\mathrm B},\varepsilon)^L S(z_{\mathrm L})a_{\mathrm L}.
  \label{eq:auxiliary-barrier-amplitude-matching}
\end{equation}
Each column contains its forward and backward amplitudes in that order;
the right origin is $L$, not zero. A fixed left interface state has unique
right amplitude coordinates if $z_{\mathrm R}\ne z_{\mathrm R}^{-1}$,
because $S(z_{\mathrm R})$ is then invertible. This does not solve the
different boundary-value problem of fixing an incident amplitude on one side
and a vanishing incoming amplitude on the other. For $L=0$ and equal lead
roots, unchanged amplitudes satisfy matching since $T^0=I_2$.

The corresponding source families are
\path{ProjectKoios/QuantumMechanics/FiniteDifference/ConstantPotential1D.lean},
\path{FiniteBarrier1D.lean} and \path{BarrierMatching1D.lean} in
\texttt{projectkoios-formal}; the last two paths use the same directory.
Their exact algebra allows complex potentials and energies. Physical
self-adjoint scattering, incoming/outgoing classification, current
normalization, transmission/reflection probabilities, continuum limits and
numerical conditioning need additional assumptions and analysis. No silicon
barrier result or probability-conservation claim follows here.

\section{Finite decisions: dominance, intervals and first passing}

Let $D$ be a nonempty finite family of candidates already admitted to the
same structural comparison problem: same spaces, source, observations,
units, gauge constraints and required channels. For each candidate $d\in D$,
let $E_i(d)\geq0$ be its error in channel $i=1,\ldots,m$, where $m\geq1$.
Channel errors may have different units; each is compared only with its own
budget $\tau_i$ in the same units. The upper-budget predicate is
$E_i(d)\leq\tau_i$ for every channel. This is not the strict lower-bound
rank gate used in silicon common-space admission.

\begin{definition}[Strict Pareto dominance]
Candidate $a$ dominates $b$ when $E_i(a)\leq E_i(b)$ for every channel and
the inequality is strict in at least one channel. A candidate is efficient
when no admitted candidate dominates it; the efficient subset is the Pareto
front. This order compares an error vector, not physical state vectors.
\end{definition}
If a dominated candidate meets all upper budgets, its dominator also meets
them. Repeating strict dominance in a finite set terminates, so any feasible
candidate has an efficient feasible weak improver. Hence a finite admitted
family is feasible exactly when its Pareto front contains a feasible member.
This does not supply a unique optimum or combine incompatible structural
gates into one numerical score. A minimizer of a weighted error sum is
efficient when every weight is strictly positive and the weighted terms have
compatible declared units: a dominator would strictly lower that sum.

Suppose instead that a channel has a justified enclosure
$\ell_i(d)\leq E_i(d)\leq u_i(d)$. The symbols $\ell_i,u_i$ denote lower
and upper endpoints, not additional tolerances. For an upper budget,
\begin{equation}
  u_i(d)\leq\tau_i\ \Longrightarrow\ \text{accepted},\qquad
  \ell_i(d)>\tau_i\ \Longrightarrow\ \text{rejected}.
  \label{eq:auxiliary-interval-decisions}
\end{equation}
If neither implication applies, that channel is unresolved. Joint acceptance
requires acceptance in every mandatory channel; one rejected channel rules
out that candidate. An unresolved candidate is not a failed candidate.
These conclusions require actual enclosures, not merely repeated binary64
agreement, and say nothing about omitted candidates or missing observations.

Let $C\subseteq\mathbb N$ be a finite ordered set of allowed cutoffs and
let $P(c)$ state that all declared gates pass at cutoff $c$. If the passing
subset is nonempty, its least member
\begin{equation}
  c_* = \min\{c\in C:P(c)\}
  \label{eq:auxiliary-first-passing-cutoff}
\end{equation}
is the unique first passing cutoff value. Its certificate must establish
$P(c_*)$ and failure at every earlier allowed cutoff. No monotonic error
law is needed for this finite statement. With interval evidence, every earlier
cutoff must be decisively rejected before the first-passing conclusion is
certified. Existence of several candidates at one cutoff does not imply a
unique candidate. Failure on a finite family is not a global impossibility
result over an unexamined continuous domain.

These constructions correspond to the source-neutral Formal families
\path{BandGauge/Certificates/FinitePareto.lean},
\path{BandGauge/Certificates/IntervalDecision.lean},
\path{BandGauge/Decision/FiniteCutoff.lean} and
\path{BandGauge/Decision/FinitePareto.lean}, beneath \path{ProjectKoios/}.
They provide decision mathematics, not campaign authority or automatic
promotion of a software diagnostic to a scientific certificate.

\section{Normalized torus averages and constrained Bessel coefficients}

Write $q=(q_1,\ldots,q_d)\in[0,1)^d$ for a dimensionless reduced torus
coordinate. The normalized average of an integrable function $f$ is
$\langle f\rangle=\int_{[0,1)^d}f(q)\,dq$, so $\langle1\rangle=1$.
For real $x$ and integer $n$, define the integer-order Bessel coefficient by
\begin{equation}
  J_n(x)=\frac1{2\pi}\int_0^{2\pi}
       e^{ix\sin\theta}e^{-in\theta}\,d\theta.
  \label{eq:auxiliary-bessel-coefficient}
\end{equation}
Here $\theta=2\pi q$ is an angle, and the Fourier coefficient uses a negative
forward phase. Expanding
$e^{ix\sin\theta}=\exp[(x/2)(e^{i\theta}-e^{-i\theta})]$
in absolutely convergent exponential series gives, for $n\geq0$,
\begin{equation}
  J_n(x)=\sum_{k=0}^{\infty}
     \frac{(-1)^k(x/2)^{n+2k}}{k!(n+k)!},\qquad
  J_{-n}(x)=(-1)^nJ_n(x).
  \label{eq:auxiliary-bessel-series}
\end{equation}
The summation index $k$ is an auxiliary nonnegative integer, not a wavevector.
Absolute convergence justifies termwise integration and extraction of the
matching exponential power. In particular
$\langle e^{ix\sin(2\pi q)}\rangle=J_0(x)$. For a sum of functions of
independent torus coordinates, Fubini factorizes the exponential average;
for example $\langle\exp[i\sum_{j=1}^d x_j\sin(2\pi q_j)]\rangle
=\prod_{j=1}^d J_0(x_j)$. This is the normalized circle/product-torus
construction, not a statement that arbitrary dependent bond phases factorize.

To expose the dependence, let $b_1,\ldots,b_m\in\mathbb Z^d$ be oriented
bond vectors and let $\mathsf B$ be the $d\times m$ integer matrix with
these columns. For the unit-modulus scalar field
$f(q)=\exp[ix\sum_{\nu=1}^m\sin(2\pi b_\nu\cdot q)]$, its coefficient
at displacement $p\in\mathbb Z^d$ is
\begin{equation}
  c_p=\sum_{\ell\in\mathbb Z^m:\,\mathsf B\ell=p}
           \prod_{\nu=1}^m J_{\ell_\nu}(x).
  \label{eq:auxiliary-constrained-bessel-coefficient}
\end{equation}
The vector $\ell$ lists Bessel orders; the integer constraint enforces the
Fourier displacement. At $p=0$ this is a constrained zero-mode sum, not in
general $J_0(x)^m$. Orientation matters because reversing a bond reverses
its sine phase.

The maintained triangular kernel uses bonds $(1,0),(0,1),(1,-1)$.
Writing $p=(p_1,p_2)$, its coefficient is
\begin{equation}
 c_p=\sum_{k\in\mathbb Z}J_{p_1-k}(x)J_{p_2+k}(x)J_k(x).
 \label{eq:auxiliary-triangular-bessel-kernel}
\end{equation}
The FCC kernel uses $e_1,e_2,e_3,e_2-e_1,e_3-e_1,e_3-e_2$, where $e_j$
are the standard integer coordinate vectors. For $p=(p_1,p_2,p_3)$ it is
\begin{equation}
 c_p=\sum_{a,b,c\in\mathbb Z}
 J_{p_1+a+b}(x)J_{p_2-a+c}(x)J_{p_3-b-c}(x)J_a(x)J_b(x)J_c(x).
 \label{eq:auxiliary-fcc-bessel-kernel}
\end{equation}
Here $a,b,c$ are auxiliary Bessel indices, unrelated to basis-change matrices
in other chapters. These are scalar Fourier constructions, not atomic
populations or authenticated electronic symmetry actions.

\subsection{Kernel truncation versus displacement-shell truncation}

Let $h=|x|/2$. The series above gives
$|J_n(x)|\leq e^{h^2}h^{|n|}/|n|!$, while the integral gives $|J_n(x)|\leq1$.
Thus an absolute bilateral sum bound is
$C(h)=e^{h^2}(2e^h-1)$. If $K\geq0$ is an integer kernel-index cap and
$h/(K+2)<1$, its omitted bilateral tail is bounded by
\begin{equation}
 T_K(h)=\frac{2e^{h^2}h^{K+1}}{(K+1)!\,[1-h/(K+2)]}.
 \label{eq:auxiliary-bessel-index-tail}
\end{equation}
Indeed the factorial-series terms after $K$ have successive ratios at most
$h/(K+2)$, so their sum is bounded by a geometric series. At $h=0$ the
omitted tail is zero. Truncating the triangular kernel to $|k|\leq K$
has coefficient remainder at most $T_K$, because its other two factors have
modulus at most one. Truncating the FCC kernel to $|a|,|b|,|c|\leq K$
has remainder at most $3T_K C(h)^2$, by a union bound on the three omitted
index directions. These are uniform coefficient bounds, not yet shell-tail
powers.

For a displacement $p$, define the dimensionless shell index
$Q(p)=p_1^2+p_2^2+p_1p_2$ in the triangular case, and
$Q(p)=p_1^2+p_2^2+p_3^2+p_1p_2+p_1p_3+p_2p_3$ in the FCC case.
Both obey $Q(p)\geq\|p\|_2^2/2$, so the shell $Q(p)\leq S$, with
$S\geq0$, is fully enumerated in the coordinate box
$|p_j|\leq\lceil\sqrt{2S}\rceil$. The shell cap $S$ and kernel cap $K$
are different truncations.

Normalized Parseval gives $\sum_p|c_p|^2=1$ because $|f|=1$.
If an approximate coefficient $\widehat c_p$ has a justified total error
bound $R_K+\eta_p$, where $R_K$ is the appropriate kernel remainder and
$\eta_p$ bounds the remaining numerical error, then
\begin{equation}
 \sum_{Q(p)>S}|c_p|^2
 \leq1-\sum_{Q(p)\leq S}
       \bigl[\max(0,|\widehat c_p|-R_K-\eta_p)\bigr]^2.
 \label{eq:auxiliary-bessel-shell-tail}
\end{equation}
The right-hand side is an exact-real conditional bound. To turn its numerical
evaluation into a rigorous enclosure, bounds must cover special-function
values, input error and the accumulation arithmetic as well as kernel
truncation. Clipping a binary64 result to $[0,1]$ does not supply that proof.

The maintained owner is
\path{python/src/ksdft2effmass/band_equivalence_hopping_locality/constrained_bessel.py}.
Its zero-mode request accepts an independent argument $x\in[0,\pi/2]$;
its tail request uses $x=2\alpha/m$, with maximum rotation amplitude
$\alpha\in[0,\pi/4]$ and bond count $m=3$ or $6$.
Geometry-specific kernel caps and the displacement-shell resource cap are
software policies, not physical convergence claims. Its SciPy/binary64 values
and explicitly supplied numerical margin are diagnostics, not a formal
interval-arithmetic certificate. The Formal normalized-circle, product-torus
and Bessel-zero families supply separate exact analysis; this appendix does
not claim they mechanize these dependent-bond kernels, runtime sums or tail
outputs.

\section{Synthetic Bain geometry and a common state space}

The selected geometry subprotocol at
\path{calculations/periodic3d_rankN_strain/geometry-common-state-space.md}
defines a synthetic homogeneous metric path, not material strain or phase
stability. Let $s\in[0,1]$ be its dimensionless path fraction and $\ell_0$
its declared synthetic direct-length unit. Set
\begin{equation}
 \begin{gathered}
 \gamma(s)=\sqrt2+s(1-\sqrt2),\quad
 t(s)=[2/\gamma(s)]^{1/3},\quad l(s)=2/t(s)^2,\\
 A_B(s)=\frac{\ell_0}{2}
 \begin{pmatrix}t&-t&t\\t&t&-t\\l&l&l\end{pmatrix}.
 \end{gathered}
 \label{eq:auxiliary-bain-metric-path}
\end{equation}
The columns of $A_B$ are primitive direct vectors in fixed Cartesian axes;
$t,l$ are dimensionless scales and $\gamma=l/t$. Direct multiplication
shows $\det A_B=\ell_0^3t^2l/2=\ell_0^3$.
The metric $G_B=A_B^{\mathsf T}A_B$ has units $\ell_0^2$, while the
reciprocal matrix $K_B=2\pi A_B^{-\mathsf T}$ has units $\ell_0^{-1}$.
A fixed reduced coordinate $q_B$ labels the same abstract torus point, but
not a fixed Cartesian wavevector $K_B(s)q_B$ along the path. The endpoints
are FCC and BCC primitive geometries, not identified phases of silicon.

Let $U$ be the exact integer coordinate bridge
\begin{equation}
 U=\begin{pmatrix}1&1&0\\0&0&1\\0&-1&0\end{pmatrix},\quad
 U^{-1}=\begin{pmatrix}1&0&1\\0&0&-1\\0&1&0\end{pmatrix},\quad \det U=1.
 \label{eq:auxiliary-bain-integer-bridge}
\end{equation}
With the endpoint orientation fixed by $A_F=A_B(0)U^{-1}$, cell coordinates
obey $n_F=Un_B$ and reciprocal coordinates obey $q_B=U^{\mathsf T}q_F$.
The integer vectors $n_B,n_F$ label cells, not internal state components.
The inverse integer matrix proves bijection on the infinite cell lattice.
For a finite quotient with positive period triple $(N_1,N_2,N_3)$,
let $D_N=\operatorname{diag}(N_1,N_2,N_3)$. A well-defined two-sided bridge
also requires preservation of $D_N\mathbb Z^3$ by $U$ and $U^{-1}$.
Integrality of $D_N^{-1}UD_N$ requires
$N_1\mid N_2$, $N_2\mid N_3$ and $N_3\mid N_2$; positivity makes this
exactly $N_2=N_3$ and $N_1\mid N_2$. Under those conditions the inverse
also preserves the period lattice. Every positive cubic quotient qualifies;
an arbitrary mixed-period quotient does not. This is applicability, not a
mesh selection or runtime-execution certificate.

For an admitted positive internal retained rank $r$, the selected common fiber is
$\mathbb C^r_{\mathrm{ret}}\oplus\mathbb C^r_{\mathrm{ext}}=\mathbb C^{2r}$,
with reference frame $F_0=[I_r\;0]^{\mathsf T}$. The infinite common space
is $\ell^2(\mathbb Z^3)\otimes\mathbb C^{2r}$, where $\ell^2(\mathbb Z^3)$
is the space of square-summable complex cell sequences. At an admitted finite
quotient its cell factor is replaced by that quotient's state space. The
endpoint map sends $|n_B,\nu\rangle$ to $|Un_B,\nu\rangle$, leaving
the internal index $\nu=0,\ldots,2r-1$ unchanged. It is a unitary
permutation, not a new physical strain response. Geometry basis offsets
classify physical support and do not add an unselected sublattice tensor
factor. For geometry endpoint labels $a,b$, let $\tau^a,\tau^b$ be the
specified dimensionless Bain-coordinate offsets; these labels are not the
internal index $\nu$. A bond with cell displacement $R\in\mathbb Z^3$
has physical vector $A_B(s)(R+\tau^b-\tau^a)$, not length determined by
its integer translation label alone.

The remaining lattice/rank-specific strain program choices stay with their
calculation-owned protocols. Geometry selection alone does not select
Hamiltonian controls, meshes, tolerances, retained evidence or execution.
This exposition neither reports those calculations as completed nor supplies
elasticity, phase-transition or material-transferability results.
